\documentclass[10pt,a4paper]{amsart}
\usepackage[margin=19mm]{geometry}
\usepackage{amsmath,amssymb,mathtools}
\usepackage[scr=rsfs,cal=euler]{mathalfa}
\usepackage{xfrac}
\usepackage[unicode,hidelinks,bookmarks=false]{hyperref}
\newcommand{\ie}{i.e.\ }
\newcommand{\cf}{cf.\ }
\newcommand{\resp}{respectively\ }
\newcommand{\Subsection}{\subsection}
\providecommand{\nobreakdash}{\nobreak\hbox{-}}
\setlength{\emergencystretch}{2em}
\multlinegap0pt
\usepackage{mathtools}
\usepackage[normalem]{ulem}
\usepackage{amssymb}
\usepackage{tikz}
\usepackage{xcolor}
\newtheorem{theorem}{Theorem}
\def\N{\mathbb{N}}
\def\Z{\mathbb{Z}}
\DeclareMathOperator{\colim}{colim}
\title{Fukaya category of symmetric product from gluing}
\author{Vivek Shende, Peng Zhou}
\date{Source: arXiv:2604.01208v1 (CC BY 4.0)}
\begin{document}
\maketitle

\noindent\emph{Source and licence.} Fukaya category of symmetric product from gluing. Version arXiv:2604.01208v1 (CC BY 4.0), distributed by arXiv under the Creative Commons Attribution 4.0 International licence, http://creativecommons.org/licenses/by/4.0/, as recorded in the arXiv licence metadata for this exact version and shown on the abstract page https://arxiv.org/abs/2604.01208v1. Full licence notice: \texttt{licenses/arxiv-2604.01208-CC-BY-4.0.txt}.\par\smallskip

\noindent\textit{Editorial scope.} Counted unit: the article's theorem pushout (source0.tex lines 1451--1455). The article's complete original proof of it is delivered with it (source0.tex lines 1472--1506, one \texttt{proof} environment of 35 lines, which the article places away from the statement). No mathematics is written, repaired or re-derived: the statement and the proof are the article's own lines, in the article's own notation, and no retained line is substituted or re-flowed.  No further source-local context is needed: every cross-reference of every retained line resolves inside the two delivered spans.  Nothing was omitted from any retained span, so the package carries no omission record.  No retained line contains a citation, so the wrapper carries no bibliography and no bibliography entry.  No retained line contains a figure environment, an image-inclusion command or drawing code, so no image is delivered and the package declares no asset dependency.  Editorial formatting only, in addition: an AMS \texttt{amsart} preamble that restates the article's own declarations and macros used by the retained lines, namely the environments \texttt{theorem} on separate counters, together with the standard packages those lines need.  The retained environments print in this document as Theorem 1 in order of appearance, whereas the article prints the same items with the numbering of its own full document.  The complete native source of the article as distributed in the arXiv e-print for this exact version is bundled unchanged as \texttt{sources/197643/source0.tex}.
\begin{theorem} \label{thm:tensor to pushout}
Let $M$ be a $\Z$-graded $\mathbf{k}[\epsilon] \otimes \mathbf{k}[\epsilon]$-module.  Then there is a canonical isomorphism
$$\left(\mathbb{T} \bigotimes_{k[\epsilon] \otimes k[\epsilon]} M\right)^n \xleftarrow{\sim} ( \mathbf{T} \otimes M^{n-1}) \!\!\!\!\! \coprod_{M^{n-1} \oplus M^{n-1}} \!\!\!\!\! M^n$$
Here, the maps $M^{n-1} \oplus M^{n-1} \to M^{n-1} \otimes \mathbf{T}$ are given by tensoring up  $\delta^0_{[2]} \oplus \delta^0_{[1]}: \mathbf{k} \oplus \mathbf{k} \to \mathbf{T}$, and the maps $M^{n-1} \oplus M^{n-1} \to M^n$ are given by the two $\epsilon$ actions.  Likewise, the $k[\epsilon]$ actions on $\mathbb{T}$ are by $\delta_{[2]}$ and $\delta_{[1]}$. 
\end{theorem}

\begin{proof}[Proof of Theorem \ref{thm:tensor to pushout}]
    Let us first consider a general expression of the form $X \otimes_A Y$, where 
    $A$ is a monoidal category and $X$ and $Y$ are, respectively, left and right $A$-modules.  Then 
    we have in general 
    $$  \colim \left( \cdots X \otimes A \otimes A \otimes Y   \,\, \substack{\rightarrow\\[-.6ex] \rightarrow \\[-.6ex] \rightarrow}   \,\, X \otimes A \otimes Y  \,\, \substack{\rightarrow\\[-.6ex] \rightarrow } \,\, X \otimes Y \right) = X \otimes_A Y.$$
    Now assume further that $A$ is $\N$-graded with $A^0 = \mathbf{k}$, $X$ is $\N$-graded, and $Y$ is $\Z$-graded.    Then, after removing redundant terms, 
    the above expression allows the degree $n$ part, $(X \otimes_A Y)^n$, to be computed as the colimit of a diagram of the following shape.   The vertices are the vertices of a barycentric subdivision of a $n$-simplex. To avoid confusion, we will refer to the vertices of the original simplex as {\em corner vertices}.  The arrows are oriented as follows: on each face, the barycenter is a source, and the corner vertices are sinks.  The barycenter of a $k$-face is labelled by a category of the form $X^{n_X} \otimes A^{m_1} \otimes \cdots \otimes A^{m_k} \otimes Y^{n_Y}$, where the superscripts are integers which sum to $n$, $n_X \ge 0$, and the $m_k > 0$;  we will indicate this category by the tuple $(n_X, m_1, m_2, \cdots, m_k, n_Y)$.  The arrows in the diagram are the monoidal structure maps, i.e. correspond to merging adjacent entries in the tuple.  The colimit is understood homotopically, i.e. over the category whose nerve is the solid simplex on which we described the diagram.  We will denote this diagram $\Delta^n(X|A|Y)$.

    We will use two basic facts about simplifying such diagrams.  If $v$ is a corner vertex, note that the sub-diagram on all vertices mapping to $v$ is precisely the star of $v$; we write $\mathrm{link}(v) := \mathrm{star}(v) \setminus v$.  There's a natural map from the colimit over $\mathrm{link}(v)$ to the value at $v$; if this map is an isomorphism, then $\Delta^n(X|A|Y)$ and $\Delta^n(X|A|Y) \setminus v$ have the same colimit of the diagram.    Second,  the colimit of $\Delta^n(X|A|Y) \setminus v$ is equivalent to the colimit on the facet of $\Delta^n(X|A|Y)$ opposing $v$. 

    Now we turn to the diagram of interest: $X = \mathbb{T}$, $A= k[\epsilon] \otimes k[\epsilon]$, and $Y = M$ arbitrary.  Observe that then $X^{\ge 3} = 0$ and $A^{\ge 3} = 0$.  Consider the corner vertex $(5,n-5)$.  The corresponding category is $\mathbb{T}^5 \otimes M^{n-5} = 0$.  The nearest neighbors are the points $(a, 5-a, n-5)$.  But every such category is zero, so indeed $(5,n-5)=0$ is the pushout of its link, and may be removed.  The same works a fortiori for $(k, n-k)$ for any $k\ge 5$, so we may iteratively collapse to the diagram on the 4-face of $\Delta^n(\mathbb{T}|\mathbf{k}[\epsilon] \otimes \mathbf{k}[\epsilon]|M)$ spanned by corner vertices $(4|n-4), \ldots, (0|n)$.  
    
    Consider the corner $(4|n-4)$.  The category there is $\mathbb{T}^4 \otimes M^{n-4}=0$.  The only non-zero immediately adjacent vertex is $(2|2|n-4)$ with value $\mathbb{T}^2  \otimes (\mathbf{k}[\epsilon] \otimes \mathbf{k}[\epsilon])^2 \otimes M^{n-2} = M^{n-2}$.  In the link of $(4|n-4)$, we have a subdiagram 
    $$0 = (3|1|n-2) \leftarrow  (2|1|1|n-2) \rightarrow  (2|2|n-2) = M^{n-2}.$$  The map 
    $$(2|1|1|n-2) = \mathbf{k} \otimes (\mathbf{k} \oplus \mathbf{k}) \otimes (\mathbf{k} \oplus \mathbf{k})  \otimes M^{n-2} \to M^{n-2}$$ is easily seen to be surjective (the $M^{n-2}$ just comes along for the ride, and the map $(2|1|1|0) \to (2|2|0)$ is obviously surjective).  It follows that the pushout of the link of $(4|n-4)$ is zero, and so we can again cut the corner and project onto the 3-simplex with corners $(3, n-3),\ldots, (0,n)$.

    Again the corner vertex vanishes: $(3|n-3) = 0$. The only potentially nonzero adjacent terms are $(2|1|n-3) = \mathbf{k} \otimes \mathbf{k}^{\oplus 2} \otimes M^{n-3}$ and $(1|2|n-3) = \mathbb{T} \otimes \mathbf{k} \otimes M^{n-3}$.  Again the $Y$ factors will come for the ride, so we suppress them (equivalently set $n=3$).   Now, 
    we have subdiagrams of the link 
    $$0 = (0|3| 0) \leftarrow (0|1|2|0) = \mathbf{k} \oplus \mathbf{k}  \twoheadrightarrow (1|2|0) = \mathbf{T} \qquad \qquad 0 = (0|3|0) \leftarrow (0|2|1|0) = \mathbf{k} \oplus \mathbf{k} \xrightarrow{\sim} (2|1|0) = \mathbf{k} \oplus \mathbf{k}. $$
    From these we see that the pushout of the link is again zero, and so we may cut the corner, and project to the 2-simplex on $(2| n-2), \ldots, (0|n)$. 

    Consider the corner $(2|n-2) = M^{n-2}$.  Its link is the diagram $(1|1|n-2) \leftarrow (0|1|1|n-2) \rightarrow (0|2|n-2)$ which we expand as
    $$\mathbf{T} \otimes (\mathbf{k} \oplus \mathbf{k}) \otimes M^{n-2} \leftarrow (\mathbf{k} \oplus \mathbf{k}) \otimes (\mathbf{k} \oplus \mathbf{k}) \otimes M^{n-2} \rightarrow \mathbf{k} \otimes  M^{n-2}$$
    We claim the pushout is in fact $M^{n-2}$.  Let's drop the $M^{n-2}$ to avoid cluttering notation, and decorate the $\mathbf{k}$ factors by $1$ and $2$ according as they come from the first or second copy of $\mathbf{k}[\epsilon]$:
    $$\mathbf{T} \otimes (\mathbf{k}_1 \oplus \mathbf{k}_2)  \leftarrow (\mathbf{k}_1 \oplus \mathbf{k}_2) \otimes (\mathbf{k}_1 \oplus \mathbf{k}_2) \rightarrow \mathbf{k} $$
    Now, the $\mathbf{k}_1 \otimes \mathbf{k}_1$ factor in the middle maps to $\mathbf{T} \otimes \mathbf{k}_1$ on the left, and dies on the right; similarly with $\mathbf{k}_2 \otimes \mathbf{k}_2$.  
    Thus the pushout of this diagram agrees with the pushout by the quotient of these pieces and their images; i.e. of 
    $$((\mathbf{T}/\mathbf{k}_1) \otimes \mathbf{k}_1) \oplus ((\mathbf{T} /\mathbf{k}_2)  \leftarrow 
    (\mathbf{k}_2 \otimes \mathbf{k}_1) \oplus (\mathbf{k}_1 \otimes \mathbf{k}_2)\rightarrow \mathbf{k}$$
    The left arrow is an isomorphism, so the colimit is $\mathbf{k}$, as desired.  We may cut the corner. 

    What remains is the diagram $(1|n-1) \leftarrow (0|1|n-1) \rightarrow (0|n)$, or, in other words
    $$\mathbf{T} \otimes M^{n-1} \leftarrow (\mathbf{k} \oplus \mathbf{k}) \otimes M^{n-1} \rightarrow M^n.$$
    This completes the proof.
\end{proof}

\end{document}
